\section{O konkretnych grupach krócej, czyli grupy Liego}

\subsection{Grupa Liego}

\begin{definition}{grupa Liego}{}
  Grupą Liego nazwiemy grupę topologiczną, która ma dodatkowo strukturę rozmaitości różniczkowej.
\end{definition}

\begin{example}[m]
  \item Grupa $U(1)$, czyli unitarnych macierzy $1\times 1$, często utożsamiana z okręgiem $S^1$.
  \item $SU(2)$ czyli znowu sferka, tym razem $S^3$.
  \item Grupa przekształceń afinicznych prostej, czyli $x\mapsto ax+b$, utożsamianymi z macierzami 
    $$\begin{bmatrix}
       a & b \\ 
      0 & 1 \end{bmatrix}$$
  \item Grupa Heisenberga, czyli macierzy $3\times 3$ górnotrójkątnych o wyrazach z $\R$,
    $$\begin{bmatrix} 1 & a  & b \\ 
      0 & 1 & c \\ 
      0 & 0 & 1 \end{bmatrix}$$
\end{example}

\subsection{Algebra Liego}

\begin{lemma}{}{}
  Lewe przesunięcie $L_g:G\to G$, $L_g(h)=gh$, indukuje izomorfizm między $T_1G$ oraz $T_gG$.
\end{lemma}

Co więcej, $TG\cong G\times T_1G$.

\begin{definition}{algebra Liego}{}
  Algebra Liego to przestrzeń styczna do grupy Liego w identyczności z mnożeniem zadanym przez nawias Liego, $[X,Y]=XY-YX$. Zwykle oznaczamy $T_1G=\mathfrak{g}$.
\end{definition}

Mimo że jeszcze nie wiemy co to takiego reprezentacja, to i tak już powiemy, że grupy Liego maja pewną fajną reprezentację.

\begin{example}
  Rozważmy automorfizm grupy Liego $G$, $A_g(h)=ghg^{-1}$. Licząc jego pochodną w identyczności dostajemy odwzorowanie $Ad_g=\frac{d}{dh}A_g(H)|_{h=1}$ algebry Liego grupy $G$. Rozważając odwzorowanie $G\ni g\mapsto Ad_g\in Aut(\mathfrak{g})$ dostajemy tzn. \acc{dołączoną reprezentacje} grupy $G$.

  Ale skoro lecimy w $T_1G$, to możemy też polecieć dalej, do przestrzeni kostycznej i dostać $Ad^*$ - reprezentację \acc{kodołączoną}.
\end{example}



